%% Glossary ------------------------------------------------------------------
\chapter{Glossary}
\label{ch:glossary}
\index{glossary}

\begin{description}[style=nextline, leftmargin=1.6em, itemsep=4pt]
  \item[ADAPT] Air Force Data Assimilative Photospheric Flux Transport model; synoptic
    magnetograms that estimate the unobserved far side of the Sun.
  \item[AIA] Atmospheric Imaging Assembly on SDO; full-disk EUV and UV imager.
  \item[Alfvén Mach number (\(M_A\))] Ratio of the shock speed to the Alfvén speed.
  \item[Alfvén speed (\(V_A\))] Speed of Alfvén waves, \(B/\sqrt{\mu_0\rho}\); depends on the
    magnetic field and plasma density.
  \item[ARTEMIS-IV] Solar radio spectrograph at Thermopylae, Greece.
  \item[Background subtraction] Removal of each frequency channel's steady receiver and sky
    level from a dynamic spectrum.
  \item[Band splitting] Division of a Type~II emission lane into two parallel bands emitted
    upstream and downstream of the shock.
  \item[Base difference] Each image minus the first image of the sequence.
  \item[\texttt{BITPIX}] FITS keyword giving the data type of the stored values.
  \item[Carrington coordinates] Heliographic coordinates that rotate with the Sun at the
    Carrington rate; used for synoptic maps.
  \item[CDAW] Coordinated Data Analysis Workshop data center at NASA GSFC, host of the
    SOHO/LASCO CME catalog.
  \item[CME] Coronal mass ejection: an eruption of magnetized plasma from the corona.
  \item[Compression ratio (\(X\))] Ratio of the downstream to the upstream plasma density at a
    shock.
  \item[Coronagraph] Telescope that blocks the solar disk with an occulter to observe the faint
    corona.
  \item[Density model] Function \(n_e(r)\) giving the coronal electron density at a
    heliocentric distance.
  \item[Derotation] Compensation for solar differential rotation over an image sequence.
  \item[Digit] Raw CALLISTO intensity value (analog-to-digital converter unit, ADU), about
    0.384~dB per digit.
  \item[Drift rate] Rate of change of the emission frequency of a burst, \(\mathrm{d}f/\mathrm{d}t\).
  \item[Dst index] Disturbance storm-time index; measure of the ring current and geomagnetic
    storm strength.
  \item[Dynamic spectrum] Radio intensity as a function of frequency and time.
  \item[e-CALLISTO] World-wide network of CALLISTO solar radio spectrometers.
  \item[EIT] Extreme-ultraviolet Imaging Telescope on SOHO.
  \item[EUV] Extreme ultraviolet.
  \item[EUVI] Extreme Ultraviolet Imager of STEREO/SECCHI.
  \item[FITS] Flexible Image Transport System, the standard astronomical data format.
  \item[Focus code] Two-digit identifier of a CALLISTO receiver or frequency band at a station.
  \item[Fold number] Factor multiplying a coronal density model to represent denser structures.
  \item[Fundamental, harmonic] Plasma emission at the local plasma frequency, or at twice it.
  \item[GCS] Graduated Cylindrical Shell, a geometric flux-rope model of CMEs.
  \item[GONG] Global Oscillation Network Group, a ground-based network providing synoptic
    magnetograms.
  \item[GOES] Geostationary Operational Environmental Satellites of NOAA, carrying the XRS,
    SGPS and SUVI instruments.
  \item[HEK] Heliophysics Event Knowledgebase.
  \item[Helioprojective coordinates] Observer-centered angular coordinates on the sky, in
    arcseconds from Sun center.
  \item[Helioviewer] Service providing JPEG2000 solar images for browsing.
  \item[HI] Heliospheric Imagers (HI1, HI2) of STEREO/SECCHI.
  \item[HMI] Helioseismic and Magnetic Imager on SDO.
  \item[J-map] Time--elongation map built from a slit through a sequence of images.
  \item[JSOC] Joint Science Operations Center, the SDO data archive.
  \item[Kp index] Planetary three-hourly index of geomagnetic activity, from 0 to 9.
  \item[LASCO] Large Angle and Spectrometric Coronagraph (C2, C3) on SOHO.
  \item[Level 1, 1.5, 1b, 2] Processing levels of solar image data, from archive products to
    registered or composite images.
  \item[NRGF] Normalizing-radial-graded filter, which flattens the radial brightness gradient of
    coronagraph images.
  \item[Occulter] Disk in a coronagraph that blocks the bright inner field of view.
  \item[Open and closed field] Magnetic field lines that reach the source surface, or return to
    the photosphere.
  \item[OriginPro style] The publication graph style used for exported figures: closed frame,
    inward ticks and a white background.
  \item[PFSS] Potential-field source-surface model of the coronal magnetic field.
  \item[Plasma frequency] Natural oscillation frequency of electrons in a plasma,
    \(f_p \approx 8.98\sqrt{n_e}\)~kHz.
  \item[Position angle (PA)] Angle on the sky measured from solar north toward east.
  \item[Power law] Function of the form \(a\,t^{-b}\).
  \item[\(R_\odot\)] Solar radius, 696\,340~km.
  \item[RFI] Radio-frequency interference from terrestrial transmitters.
  \item[Ridge] The line of peak intensity of a burst in a dynamic spectrum.
  \item[Running difference] Each image minus the previous image.
  \item[SDO] Solar Dynamics Observatory.
  \item[SECCHI] Sun Earth Connection Coronal and Heliospheric Investigation on STEREO (EUVI,
    COR1, COR2, HI1, HI2).
  \item[SEP] Solar energetic particles.
  \item[SGPS] Solar and Galactic Proton Sensor on GOES-R series satellites.
  \item[SOHO] Solar and Heliospheric Observatory.
  \item[Source surface] Sphere in the PFSS model beyond which the field is radial.
  \item[STEREO] Solar Terrestrial Relations Observatory (Ahead and Behind spacecraft).
  \item[Stonyhurst coordinates] Heliographic coordinates with longitude zero facing the Earth.
  \item[SUVI] Solar Ultraviolet Imager on GOES.
  \item[SWAVES] Radio and plasma wave instrument on STEREO.
  \item[Synoptic map] Full-Sun map assembled over a solar rotation.
  \item[Type II burst] Slowly drifting radio burst produced by a shock wave.
  \item[Type III burst] Fast-drifting radio burst produced by electron beams.
  \item[UT] Universal Time.
  \item[VSO] Virtual Solar Observatory.
  \item[XRS] X-Ray Sensor on GOES, with short (0.05--0.4~nm) and long (0.1--0.8~nm) channels.
\end{description}
